\newpage
\section{Acta Número 18}
\label{sec:acta18}

\noindent \textbf{Fecha de la sesión:}\par
Lunes, 4 de Mayo de 2026

\vspace*{0.5cm}
\noindent \textbf{Datos de la Sesión:}
\vspace{-0.2cm}
\begin{itemize}[noitemsep]
    \item \textbf{Modalidad:} Virtual - Google Meet.
    \item \textbf{Hora de Inicio:} 18:00
    \item \textbf{Hora de Fin:} 20:00
\end{itemize}

\vspace*{0.5cm}
\noindent \textbf{Orden del Día:}
\vspace{-0.2cm}
\begin{itemize}[noitemsep]
    \item Control de Asistencia.
    \item Evento de \textit{Sprint Planning} para el Sprint 3.
    \item Inicio de fase de Aseguramiento de Calidad (QA) para los entregables de los Sprints 1 y 2.
    \item Gestión de errores y ejecución de correcciones (Sprints 1 y 2).
    \item Firmas y Aprobación de los Presentes.
\end{itemize}

\vspace*{0.5cm}
\noindent \textbf{Socios Fundadores Presentes:}
\vspace{-0.2cm}
\begin{enumerate}[noitemsep]
    \item Pardo Pozo Juan Diego (Jefe de Proyecto).
    \item Arnez Jaimes Mauricio.
    \item Quispe Flores Camila Belen.
    \item Ariñez Bautista Jim Gabriel.
    \item Medina Rodríguez Mateo Jhoustin.
    \item Molina Maldonado Tomas Manuel.
    \item Gutierrez Guzmán Angheline.
\end{enumerate}

\newpage
\subsection{Desarrollo del Acta}

\noindent \textbf{Control de Asistencia del Team}\par
Se procedió a la toma de asistencia a la hora acordada, corroborando la presencia de los 7 socios fundadores de la grupo-empresa Byte Busters S.R.L. Al contar con el quórum necesario, se dio inicio formal a la sesión de trabajo.

\vspace*{0.5cm}
\noindent \textbf{Evento de \textit{Sprint Planning} - Sprint 3}\par
Con la participación y el acuerdo de todo el equipo de desarrollo, se llevó a cabo la ceremonia de \textit{Sprint Planning} correspondiente al tercer sprint. Inicialmente, se estableció de manera clara y compartida el \textit{Sprint Goal}, alineando los esfuerzos técnicos con el valor de negocio esperado. Posteriormente, se realizó la selección estratégica de los \textit{Product Backlog Items} (PBIs) más prioritarios, trasladándolos al \textit{Sprint Backlog}. Para asegurar la viabilidad del ciclo, el equipo calculó su capacidad técnica real (\textit{Capacity}) y reafirmó su compromiso con el trabajo seleccionado. Finalmente, se revisó y validó en consenso la Definición de Hecho (\textit{Definition of Done} - DoD) aplicable a todas las tareas del sprint.

\vspace*{0.5cm}
\noindent \textbf{Inicio de Aseguramiento de Calidad (QA) - Sprints 1 y 2}\par
De forma paralela a la planificación de los nuevos desarrollos, se determinó por unanimidad el inicio formal de los protocolos de Aseguramiento de Calidad (QA) sobre el incremento funcional de los Sprints 1 y 2. Se definieron los escenarios de prueba a ejecutar, garantizando que todos los flujos de la arquitectura base sean sometidos a rigorosas validaciones de integración y funcionamiento, asegurando que el producto cumpla con los estándares de calidad definidos por la empresa antes de su despliegue final.

\vspace*{0.5cm}
\noindent \textbf{Gestión de Errores y Correcciones - Sprints 1 y 2}\par
Alineados con la fase de QA, el equipo estableció el flujo operativo para la gestión de errores encontrados en los entregables previos (Sprints 1 y 2). Se acordó un proceso estructurado para la detección, registro (\textit{triage}), análisis de causa raíz y rápida resolución de bugs funcionales y técnicos. Todos los desarrolladores confirmaron su compromiso para aplicar los \textit{fixes} necesarios en tiempo y forma, asegurando la eliminación definitiva de la deuda técnica pendiente y la estabilización completa del sistema.

\vspace*{0.5cm}
\noindent \textbf{Aprobación Oficial}\par
Habiendo concluido exitosamente la planificación del tercer sprint, con metas, capacidad y métricas de calidad claramente definidas y acordadas por todos, el equipo expresa su total conformidad. Sin más puntos técnicos ni estratégicos que tratar, se da por aprobada la sesión y se procede con la firma correspondiente.

\newpage
\noindent \textbf{Firmas y Aprobación de los Presentes}
\begin{center}
    \noindent
    \setlength{\tabcolsep}{0pt}
    \begin{tabular}{ccc}
        \espacioFirma{assets/img/firmas/mauricio.jpg}{Mauricio Jaimes Arnez}{13751221} & 
        \espacioFirma{assets/img/firmas/camila.jpg}{Camila Belen Quispe Flores}{13097244} &
        \espacioFirma{assets/img/firmas/jim.jpg}{Jim Gabriel Ariñez Bautista}{9517642} \\[1.0cm]
        \espacioFirma{assets/img/firmas/mateo.jpg}{Mateo Medina Rodríguez}{9453809} &
        \espacioFirma{assets/img/firmas/tomas.jpg}{Tomas Molina Maldonado}{8051701} & 
        \espacioFirma{assets/img/firmas/angheline.jpg}{Angheline Gutierrez Guzmán}{13751815} \\[1.0cm]
    \end{tabular}
    \vspace{0.2cm} 
    \begin{tabular}{cc}
        \espacioFirma{assets/img/firmas/diego.jpg}{Juan Diego Pardo Pozo}{12747783}
    \end{tabular}
\end{center}

\vspace{0.5cm}
\noindent \textbf{Fecha de última revisión:} Lunes, 4 de Mayo de 2026